\vspace{-0.7em}
\section{Knowledge-Aware Refusal in Factual Tasks}
\vspace{-0.7em}
\label{sec:bg}

\noindent\textbf{Knowledge-Aware Refusal.}
Knowledge-aware refusal measures whether a model can appropriately decline to answer questions it doesn't know.
When we define ``knowing'' as the ability to provide correct answers, knowledge-aware refusal capability can be quantified by the alignment between a model's refusal decisions and its answer incorrectness.
Specifically, a model good at knowledge-aware refusal shows low refusal rates for questions it can answer correctly, while showing high refusal rates for questions it cannot answer accurately.
This knowledge-aware refusal ability is fundamental for reliable deployment of LLMs in factual tasks. 
Previous works have explored to evaluate this capability~\citep{cheng2024aiassistantsknowdont, kapoor2024large}. 
Among existing approaches, two types of metrics, refusal-rate-based and calibration-based metrics, have been employed to measure knowledge-aware refusal. 
However, both exhibit distinct limitations in accurately assessing this ability.

\vspace{-0.5em}
\begin{table}[thbp]
    \centering
    \footnotesize
    \caption{Baseline factuality metrics used for comparison. $c$ and $r$ denote correct answer rate and refusal rate, respectively.}
    \vspace{-0.5em}
    \begin{tabularx}{\textwidth}{l l l}
        \toprule
        \textbf{Metric} & \textbf{Formula} & \textbf{Definition} \\
        \midrule
        Correct Answer Rate & $c$ & Proportion of correct answers \\
        Refusal Rate & $r$ & Proportion of refusal answers \\
        Correct given Attempted (C/A) & $c/(1 - r)$ & Correct answer rate among answered questions \\
        F-score & $2c/(2 - r)$ & Harmonic mean of Correct Answer Rate and C/A \\
        Weighted Score & $c - p(1-r)$ & Weighted difference of $c$ and $r$ \\
        \midrule
        Refusal Index & Eq. (\ref{eq:ri-estimation}) & Correlation between refusal and answer incorrectness \\
        \bottomrule
    \end{tabularx}
    \label{tab:baseline_metrics}
    \vspace{-0.5em}
\end{table}


\noindent\textbf{Limitations of Refusal-Rate-Based Metrics.}
Refusal rate alone only measures the frequency of refusals, without capturing the correlation between refusal and answer correctness.
For instance, one can prompt an LLM to be more cautious, thereby increasing the refusal rate without actually improving the model's knowledge-aware refusal ability.
To address this limitation, recent works have combined refusal rates with correct answer rates for evaluation.
The underlying intuition is that if a model refuses more samples while maintaining its correct answer rate, it demonstrates a stronger ability to identify uncertain questions.
For example, SimpleQA~\citep{wei2024measuringshortformfactualitylarge} employs an F1 score between the correct answer rate within answered questions and the refusal rate to balance over-refusal and over-confidence. 
We list these combined metrics in Table~\ref{tab:baseline_metrics}.
However, such combinations of refusal rates and correct answer rates are heuristics designed to penalize over-refusal, which fail to capture the fundamental correlation between refusal and incorrectness. 
As our experiments demonstrate in \Secref{sec:refusal_index_validation}, these metrics do not measure a consistent underlying ability: when prompting models to increase or decrease their refusal rates, these metric values fluctuate significantly (e.g. F1 score used in SimpleQA varies by up to 70\%)

\noindent\textbf{Limitations of Calibration-Based Metrics}
Other works employ calibration methods, which first estimate the uncertainty (or conversely, confidence score) of model outputs, then compute the correlation between confidence scores and answer correctness.
Metrics such as Expected Calibration Error (ECE) and AUROC are commonly used to quantify this relationship.
To apply calibration methods, previous works have designed various approaches to estimate the uncertainty of model outputs. 
For example, some methods~\citep{xiong2023can} instruct models to assign verbalized confidence scores to their own outputs, while others~\citep{ulmer-etal-2024-calibrating} train auxiliary models to predict confidence from text outputs. 
A more faithful approach involves repeatedly sampling multiple outputs for a single question and using the frequency of refusal answers as an uncertainty measure.
However, the confidence scores derived from these methods cannot fully represent a model's refusal probability. 
First, studies eliciting verbalized confidence report systematic overconfidence and high ECE, while asking the same model to vote across samples (e.g., SimpleQA) yields frequency-based curves much closer to the diagonal for larger models~\citep{xiong2023can,wei2024measuringshortformfactualitylarge}. 
Second, auxiliary calibrators such as APRICOT or rank-calibration frameworks can produce near-perfect ECE/AUROC~\citep{ulmer-etal-2024-calibrating,huang-etal-2024-rank}, yet these numbers mainly reflect the auxiliary predictor or ranking procedure rather than the base model's refusal behavior. 
Third, white-box confidence proxies like P(True) can even appear well calibrated on multiple-choice settings~\citep{kadavath2022mostly}, further showing that calibration verdicts swing with the chosen estimator. 
Therefore, while these methods provide valuable insights into calibration properties within LLMs, their uncertainty estimates do not directly reflect model refusal probability. 
In real-world applications, we expect models to abstain from providing uncertain answers. 
Thus, measuring the correlation between refusal behavior and incorrectness in black-box settings remains an unsolved yet crucial challenge for assessing the factual reliability of language models. Appendix~\ref{app:external-calibration} provides an empirical comparison of three representative calibrators ($P(\mathrm{IK})$, APRICOT, and $P(\mathrm{Answering})$) on Qwen3-32B, showing that they disagree and that only the sampling-based method exposes the over-confidence captured by RI.

\begin{remark}[Properties of Effective Measurement]
We identify three key properties that an effective metric for knowledge-aware refusal should satisfy:
\begin{enumerate}[leftmargin=*]
    \item \textbf{Faithful:} Accurately quantify knowledge-aware refusal capability.
    \item \textbf{Consistent:} Remain stable across different refusal rates induced by varying prompts or instructions.
    \item \textbf{Direct:} Derive directly from black-box LLM refusal decisions, without relying on auxiliary models.
\end{enumerate} 
\end{remark}
